\chapter{Методика}
\label{ch:methods}

%надо провести адиабатический процесс, чтобы была связь с помощью уравнения Пуассона. Есть сосуд фиксированного объёма на основе этого создать адиабатический процесс, на основе этого описать методику
%входное выходное давление. как сделать так, чтобы температуру не учитывать, вводим изотермический процесс

Для того, чтобы измерить скорость откачки вакуума из сосуда некоторого объёма при помощи диффузионного насоса, можно провести процесс откачки остаточного газа от некоторого начального давления до предельного. Можно получить следующее выражение для зависимости давления от времени в процессе откачки (вывод см. в Приложении А):

\begin{equation}
    P-P_\text{пр}=(P_0 - P_\text{пр})\exp\left(-\frac{W}{V}t\right),
\end{equation}

где $P$ - текущее давление в сосуде, $P_0$ - начальное давление, $P_\text{пр}$ - предельное давление, $V$ - откачиваемый объём, $W$ - скорость откачки, $t$ - время, прошедшее с начала откачки. Величина $\tau=V/W$ в этом выражении является мерой эффективности откачной системы.

Следовательно, для определения скорости откачки по улучшению вакуума необходимо знать объём сосуда $V$, начальное давление $P_0$ и зависимость давления от времени $P(t)$ в процессе откачки. Давление в сосуде можно измерить, например, при помощи ионизационного манометра. Объём можно найти, например, измерив давление некоторого известного объёма воздуха при известном начальном давлении после заполнения им всего объёма установки и использовав закон Бойля - Мариотта .


